\subsection{g-模的张量积}

设 \(\mathrm{E},\mathrm{F}\) 为 \(\mathfrak{g}\)-模。对于所有 \(\lambda, \mu \in \mathfrak{h}^*\)，有 \(\mathrm{E}^\lambda \otimes \mathrm{F}^\mu \subset (\mathrm{E} \otimes \mathrm{F})^{\lambda + \mu}\)（第 VII 章，§1，第 1 节，命题 2 (ii)）。若 \(\mathrm{E}\) 和 \(\mathrm{F}\) 是有限维的，则 \(\mathrm{E} = \sum_{\lambda \in \mathrm{P}} \mathrm{E}^\lambda\) 且 \(\mathrm{F} = \sum_{\mu \in \mathrm{P}} \mathrm{F}^\mu\)；因此，

\[
(\mathrm{E} \otimes \mathrm{F}) ^ {\nu} = \sum_ {\lambda , \mu \in \mathrm{P}, \lambda + \mu = \nu} \mathrm{E} ^ {\lambda} \otimes \mathrm{F} ^ {\mu}.
\]

换言之，赋予 \(E \otimes F\) 其 P 型分次后，它是分次向量空间 E 和 F 的分次张量积。

命题 9. 设 E、F 是有限维单 g-模，最高权分别为 \(\lambda,\mu\)。

\begin{enumerate}
\def\labelenumi{(\roman{enumi})}
\item
  \(\mathrm{E} \otimes \mathrm{F}\) 中最高权为 \(\lambda + \mu\) 的分量是一个单子模，由 \((\mathrm{E} \otimes \mathrm{F})^{\lambda + \mu} = \mathrm{E}^{\lambda} \otimes \mathrm{F}^{\mu}\) 生成。
\item
  \(\mathrm{E} \otimes \mathrm{F}\) 的任意单子模的最高权都 \(\leq \lambda + \mu\)（参见 §9，命题 2）。
\end{enumerate}

若 \(\alpha, \beta \in \mathrm{P}\) 且 \(\mathrm{E}^{\alpha} \otimes \mathrm{F}^{\beta} \neq 0\)，则 \(\alpha \leq \lambda\) 且 \(\beta \leq \mu\)。因此，\((\mathrm{E} \otimes \mathrm{F})^{\lambda + \mu}\) 等于 \(\mathrm{E}^{\lambda} \otimes \mathrm{F}^{\mu}\)，从而维数为 1，而 \(\lambda + \mu\) 是 \(\mathrm{E} \otimes \mathrm{F}\) 的最高权。 \(\mathrm{E}^{\lambda} \otimes \mathrm{F}^{\mu}\) 的每个非零元素都是本原元。根据 §6 第 2 节的命题 4，\(\mathrm{E} \otimes \mathrm{F}\) 中最高权为 \(\lambda + \mu\) 的同构型分支的长度为 1。

备注。保留命题 9 的记号。设 C 是 \(E \otimes F\) 中最高权为 \(\lambda + \mu\) 的同构型分支。则 C 只依赖于 E 和 F，而不依赖于 h 和基 B 的选择。换言之，设 \(h'\) 是 g 的一个分裂嘉当子代数，\(R'\) 是 \((\mathfrak{g}, \mathfrak{h}')\) 的根系，\(B'\) 是 \(R'\) 的一个基；设 \(\lambda', \mu'\) 是 E、F 关于 \(h'\) 和 \(B'\) 的最高权；设 \(C'\) 是 \(E \otimes F\) 中最高权为 \(\lambda' + \mu'\) 的同构型分支；则 \(C' = C\)。事实上，为证明这一点，可以通过扩张基域假设 k 是代数闭域。于是存在 \(s \in \operatorname{Aut}_{e}(\mathfrak{g})\)，将 \(h\) 映到 \(h'\)，将 R 映到 \(R'\)，将 B 映到 \(B'\)。设 \(S \in \mathbf{SL}(E \otimes F)\) 具有第 1 节命题 2 中的性质。于是 \(\mathrm{S}((\mathrm{E} \otimes \mathrm{F})^{\lambda + \mu}) = (\mathrm{E} \otimes \mathrm{F})^{\lambda' + \mu'}\)，且 \(\mathrm{S}(\mathrm{C}) = \mathrm{C}\)。因此 \((\mathrm{E} \otimes \mathrm{F})^{\lambda' + \mu'} \subset \mathrm{C}' \cap \mathrm{S}(\mathrm{C}) = \mathrm{C}' \cap \mathrm{C}\)，所以 \(C' = C\)。这样，两个有限维单 g-模的同构类可典则地确定第三个同构类；换言之，我们在有限维单 g-模的同构类集合 \(S_{g}\) 上定义了一个复合律。在此结构下，\(S_{g}\) 典则同构于加法幺半群 \(P_{++}\)。

推论 1. 设 \((\varpi_{\alpha})_{\alpha \in \mathrm{B}}\) 是关于 B 的基本权族。设 \(\lambda = \sum_{\alpha \in \mathrm{B}} m_{\alpha} \varpi_{\alpha} \in \mathrm{P}_{++}\)。对于所有 \(\alpha \in \mathrm{B}\)，设 \(\mathrm{E}_{\alpha}\) 是单 \(\mathfrak{g}\)-模，其最高权为 \(\varpi_{\alpha}\)。在 \(\mathfrak{g}\)-模 \(\bigotimes_{\alpha \in \mathrm{B}} (\bigotimes^{m_{\alpha}} \mathrm{E}_{\alpha})\) 中，最高权为 \(\lambda\) 的同构型分支长度为 1。

这由命题 9 对 \(\sum_{\alpha \in \mathrm{B}}m_{\alpha}\) 作归纳得到。

推论 2. 假设 k 是 R、C 或非离散完备超度量域。设 G 是李代数为 g 的李群。假设对于 g 的任意基本表示 \(\rho\)，都存在 G 的一个解析线性表示 \(\rho'\)，使得 \(\rho = \mathrm{L}(\rho')\)。那么对于 g 的任意有限维线性表示 \(\pi\)，都存在 G 的一个解析线性表示 \(\pi'\)，使得 \(\pi = \mathrm{L}(\pi')\)。

使用推论 1 的记号。存在一个 G 的表示 \(\sigma\)，作用在 \(X = \bigotimes_{\alpha \in B} (\bigotimes^{m_{\alpha}} E_{\alpha})\) 上，使得 \(L(\sigma)\) 对应于 X 的 \(\mathfrak{g}\)-模结构（第 III 章，§3，第 11 节，命题 41 的推论 3）。设 C 是 X 中最高权为 \(\lambda\) 的同构型分支。根据第 III 章，§3，第 11 节，命题 40，只需证明 C 在 \(\sigma(G)\) 下稳定。设 \(g \in G\)，且 \(\varphi = \operatorname{Ad}(g)\)。则 \(\sigma(g)a_X \sigma(g)^{-1} = (\varphi(a))_X\) 对所有 \(a \in \mathfrak{g}\) 成立。另一方面，\(\varphi\) 是 \(\mathfrak{g}\) 的一个自同构，将 \(\mathfrak{h}\) 映到 \(\mathfrak{h}'\)，将 R 映到 \(R' = R(\mathfrak{g}, \mathfrak{h}')\)，将 B 映到 \(B'\)，它是 \(R'\) 的一个基，并将 \(\varpi_{\alpha}\) 映到 \(\varpi_{\alpha}'\)，即 \(E_{\alpha}\) 关于 \(\mathfrak{h}'\) 和 \(B'\) 的最高权（因为 \(\varphi\) 将 \(E_{\alpha}\) 变为一个 \(\mathfrak{g}\)-模，该模与 \(E_{\alpha}\) 同构）。因此 \(\varphi\) 将 \(\lambda\) 映到 \(\sum m_{\alpha} \varpi_{\alpha}'\)。根据上面的备注，\(\sigma(g)(C) = C\)。

命题 10. 设 \(\lambda, \mu \in \mathrm{P}_{++}\)。设 E、F、G 是单 \(\mathfrak{g}\)-模，其最高权分别为 \(\lambda, \mu, \lambda + \mu\)。设 \(\mathcal{X}\)（分别为 \(\mathcal{X}', \mathcal{X}''\)）是 E（分别为 F、G）的权集。则 \(\mathcal{X}'' = \mathcal{X} + \mathcal{X}'\)。

我们有 \(E = \bigoplus_{\nu \in P} E^{\nu}, F = \bigoplus_{\sigma \in P} F^{\sigma}\)，所以 \(E \otimes F\) 是下列各空间的直和

\[
(\mathrm{E} \otimes \mathrm{F}) ^ {\tau} = \sum_ {\nu + \sigma = \tau} \mathrm{E} ^ {\nu} \otimes \mathrm{F} ^ {\sigma}.
\]

根据命题 9，\(\mathbf{G}\) 可识别为一个 \(\mathfrak{g}\)-子模，并且它包含在 \(\mathrm{E} \otimes \mathrm{F}\) 中，因此 \(\mathcal{X}'' \subset \mathcal{X} + \mathcal{X}'\)。我们有 \(\mathrm{G}^{\tau} = \mathrm{G} \cap (\mathrm{E} \otimes \mathrm{F})^{\tau}\)，只需证明：对于 \(\nu \in \mathcal{X}\) 和 \(\sigma \in \mathcal{X}'\)，有 \(\mathrm{G} \cap (\mathrm{E} \otimes \mathrm{F})^{\nu + \sigma} \neq 0\)。设 \((e_1, \ldots, e_n)\)（分别为 \((f_1, \ldots, f_p))\) 是 \(\mathrm{E}\)（分别为 \(\mathrm{F}\)）的一个基，其中每个基向量属于某个 \(\mathrm{E}^\nu\)（分别为 \(\mathrm{F}^\sigma\)），并且 \(e_1 \in \mathrm{E}^\lambda\)（分别为 \(f_1 \in \mathrm{F}^\mu\)）。\(e_i \otimes f_j\) 构成 \(\mathrm{E} \otimes \mathrm{F}\) 的一个基。假设要证明的结论不成立。则存在一对 \((i,j)\)，使得每个 \((i,j)\) 号坐标在 \(\mathrm{G}\) 的每个元素中都为零。设 \(\mathrm{U}\) 是 \(\mathfrak{g}\) 的包络代数，\(\mathrm{U}'\) 是 \(\mathrm{U}\) 的对偶，\(c\) 是 \(\mathrm{U}\) 的余积。对于所有 \(u \in \mathrm{U}\)，令 \(x_i(u)\)（分别为 \(y_j(u)\)）为 \(u(e_1)\)（分别为 \(u(f_1)\)）的 \(i\)（分别为 \(j\)）号坐标；令 \(z_{ij}(u)\) 为编号为 \((i,j)\) 的 \(u(e_1 \otimes f_1)\) 的坐标。于是 \(x_i, y_j, z_{ij} \in \mathrm{U}'\)。现在 \(e_1\) 生成 \(\mathfrak{g}\)-模 \(\mathrm{E}\)，所以 \(x_i \neq 0\)，同样 \(y_j \neq 0\)。根据 \(\mathfrak{g}\)-模 \(\mathrm{E} \otimes \mathrm{F}\) 的定义（第 I 章，§3，第 2 节），若 \(c(u) = \sum u_s \otimes u_s'\)，则有

\[
z _ {i j} (u) = \sum_ {s} x _ {i} (u _ {s}). y _ {j} (u _ {s} ^ {\prime}) = \langle c (u), x _ {i} \otimes y _ {j} \rangle .
\]

换言之，\(z_{ij}\) 是 \(x_{i}\) 与 \(y_{j}\) 在代数 \(U'\) 中的乘积。但这个代数是整环（第 II 章，§1，第 5 节，命题 10），所以 \(z_{ij} \neq 0\)。由于 \(u(e_{1} \otimes f_{1}) \in \mathrm{G}\) 对所有 \(u \in U\) 成立，这就矛盾了。

